\documentclass[11pt]{article}
\usepackage[margin=1in]{geometry}
\usepackage[T1]{fontenc}
\usepackage{lmodern,microtype,hyperref,enumitem,amsmath,amssymb,xcolor,booktabs}
\usepackage[numbers,sort&compress]{natbib}
\hypersetup{colorlinks=true,urlcolor=blue,citecolor=blue,linkcolor=blue}
\title{Q30: Comprehension Check}
\author{DLS 2026 --- Kukla}
\date{October 2026}
\begin{document}
\maketitle
\textbf{Comprehension check.} During training of a deep network you hit out-of-memory, but the \emph{model
parameters} comfortably fit in GPU memory. What is most likely consuming the memory, and name a technique
that reduces it by doing extra computation.

\vspace{0.8cm}\hrule\vspace{4pt}
{\small\textbf{Answer.} The \textbf{activations}---the intermediate forward-pass tensors that reverse-mode
autodiff must keep alive to compute gradients in the backward pass---are most likely consuming the memory.
Activation memory scales with network depth and batch size and typically dominates parameter memory during
training. \textbf{Gradient checkpointing (rematerialization)} reduces it: instead of storing all
activations, you save only a subset and \emph{recompute} the others during the backward pass, trading a
second partial forward (extra compute) for a large reduction in peak memory.}
\end{document}
